\documentclass[11pt,a4paper]{article}

% --- Codificación y tipografía ---
\usepackage[utf8]{inputenc}
\usepackage[T1]{fontenc}
\usepackage{lmodern}

% --- Diseño de página y tablas ---
\usepackage[margin=2.3cm]{geometry}
\usepackage{array}
\usepackage{longtable}
\usepackage[table]{xcolor}
\usepackage{titlesec}
\usepackage{hyperref}

\hypersetup{
  colorlinks=true,
  linkcolor=teal!60!black,
  urlcolor=teal!60!black,
  pdftitle={Esquema de tablas en Cassandra - Papaws}
}

% Colores corporativos aproximados del proyecto
\definecolor{moss}{HTML}{005F73}
\definecolor{mosslight}{HTML}{D7E8E6}
\definecolor{clay}{HTML}{BB3E03}

\titleformat{\section}{\Large\bfseries\color{moss}}{\thesection}{0.6em}{}
\titleformat{\subsection}{\large\bfseries\color{clay}}{\thesubsection}{0.6em}{}

% Tipos de columna a ancho fijo
\newcolumntype{C}{>{\ttfamily}p{4.6cm}}   % nombre de columna (monoespaciado)
\newcolumntype{T}{>{\ttfamily}p{3.2cm}}   % tipo de dato (monoespaciado)
\newcolumntype{R}{p{6.0cm}}               % rol / descripción

% Encabezado de tabla reutilizable
\newcommand{\cabecera}{%
  \rowcolor{moss}\textcolor{white}{\textbf{Columna}} &
  \textcolor{white}{\textbf{Tipo}} &
  \textcolor{white}{\textbf{Rol}} \\
}

\title{\textbf{\color{moss}Esquema de tablas en Cassandra}\\[0.3em]
\large Proyecto Papaws}
\author{}
\date{\today}

\begin{document}

\maketitle

\noindent
El sistema usa \textbf{dos keyspaces} independientes, uno por microservicio, ambos
con la estrategia \texttt{SimpleStrategy} y \texttt{replication\_factor = 2}. El
modelo sigue el patrón de Cassandra de \emph{una tabla por consulta} (\emph{query-first}):
los datos se desnormalizan y se replican en varias tablas para poder leerlos por
distintas claves sin usar \texttt{ALLOW FILTERING}.

\vspace{0.4em}
\noindent\textbf{Leyenda:}\;
\colorbox{mosslight}{\textbf{PP}} = clave de partición \quad
\colorbox{mosslight}{\textbf{CC}} = columna de \emph{clustering} \quad
\textemdash{} = columna regular.

% ============================================================
\section{Keyspace \texttt{papaws\_animales}}
% ============================================================
Servicio de Rescatistas y Animales.

% ----- rescatistas_by_id -----
\subsection{\texttt{rescatistas\_by\_id}}
Catálogo de rescatistas. Incluye el rescatista interno \emph{Refugio}
(\texttt{oculto = true}), destino por defecto al reasignar animales.

\begin{longtable}{C T R}
\cabecera
\endfirsthead
\cabecera
\endhead
\texttt{id} & uuid & \textbf{PP} \textemdash{} identificador del rescatista \\
nombre\_completo & text & \textemdash \\
telefono\_contacto & text & \textemdash \\
correo\_electronico & text & \textemdash \\
organizacion & text & \textemdash \\
zona\_operacion & text & \textemdash \\
activo & boolean & Borrado lógico (\texttt{false} = dado de baja) \\
oculto & boolean & No se muestra en la gestión (p.\,ej. Refugio) \\
\end{longtable}

% ----- animales_by_id -----
\subsection{\texttt{animales\_by\_id}}
Tabla principal de animales, consultada por su identificador.

\begin{longtable}{C T R}
\cabecera
\endfirsthead
\cabecera
\endhead
\texttt{id} & uuid & \textbf{PP} \textemdash{} identificador del animal \\
nombre & text & \textemdash \\
especie & text & \textemdash \\
peso\_actual & decimal & \textemdash \\
fecha\_ingreso & timestamp & \textemdash \\
rescatista\_id & uuid & Referencia al rescatista \\
\end{longtable}

% ----- animales_by_rescatista -----
\subsection{\texttt{animales\_by\_rescatista}}
Réplica desnormalizada para listar los animales de un rescatista
(usada en la reasignación al eliminar un rescatista).

\begin{longtable}{C T R}
\cabecera
\endfirsthead
\cabecera
\endhead
\texttt{rescatista\_id} & uuid & \textbf{PP} \textemdash{} agrupa por rescatista \\
\texttt{id} & uuid & \textbf{CC} \textemdash{} identificador del animal \\
nombre & text & \textemdash \\
especie & text & \textemdash \\
peso\_actual & decimal & \textemdash \\
fecha\_ingreso & timestamp & \textemdash \\
\end{longtable}

% ============================================================
\section{Keyspace \texttt{papaws\_consulta}}
% ============================================================
Servicio de Consultas, Veterinarios, Servicios y Productos.

% ----- veterinarios_by_id -----
\subsection{\texttt{veterinarios\_by\_id}}
\begin{longtable}{C T R}
\cabecera
\endfirsthead
\cabecera
\endhead
\texttt{id} & uuid & \textbf{PP} \textemdash{} identificador del veterinario \\
nombre\_completo & text & \textemdash \\
telefono\_contacto & text & \textemdash \\
especialidad\_principal & text & \textemdash \\
activo & boolean & Borrado lógico \\
\end{longtable}

% ----- servicios_by_id -----
\subsection{\texttt{servicios\_by\_id}}
\begin{longtable}{C T R}
\cabecera
\endfirsthead
\cabecera
\endhead
\texttt{id} & uuid & \textbf{PP} \textemdash{} identificador del servicio \\
nombre & text & \textemdash \\
descripcion & text & \textemdash \\
duracion\_estimadaminutos & int & Duración estimada (minutos) \\
precio\_base & decimal & \textemdash \\
activo & boolean & Borrado lógico \\
\end{longtable}

% ----- productos_by_id -----
\subsection{\texttt{productos\_by\_id}}
\begin{longtable}{C T R}
\cabecera
\endfirsthead
\cabecera
\endhead
\texttt{id} & uuid & \textbf{PP} \textemdash{} identificador del producto \\
nombre & text & \textemdash \\
tipo & text & \textemdash \\
unidad\_medida & text & \textemdash \\
stock\_disponible & int & Inventario disponible \\
activo & boolean & Borrado lógico \\
\end{longtable}

% ----- consultas_by_codigo -----
\subsection{\texttt{consultas\_by\_codigo}}
Tabla principal de consultas, consultadas por su código.

\begin{longtable}{C T R}
\cabecera
\endfirsthead
\cabecera
\endhead
\texttt{codigo} & text & \textbf{PP} \textemdash{} código de la consulta \\
fecha\_hora & timestamp & \textemdash \\
estado & text & Pendiente / Confirmada / Completada / Cancelada \\
observaciones & text & \textemdash \\
diagnostico & text & \textemdash \\
indicaciones\_seguimiento & text & \textemdash \\
animal\_id & uuid & Referencia al animal \\
veterinario\_id & uuid & Referencia al veterinario \\
\end{longtable}

% ----- consulta_servicios_by_codigo -----
\subsection{\texttt{consulta\_servicios\_by\_codigo}}
Servicios asociados a cada consulta (relación N a N).

\begin{longtable}{C T R}
\cabecera
\endfirsthead
\cabecera
\endhead
\texttt{codigo} & text & \textbf{PP} \textemdash{} código de la consulta \\
\texttt{servicio\_id} & uuid & \textbf{CC} \textemdash{} servicio asociado \\
\end{longtable}

% ----- consulta_productos_by_codigo -----
\subsection{\texttt{consulta\_productos\_by\_codigo}}
Productos utilizados en cada consulta y la cantidad consumida.

\begin{longtable}{C T R}
\cabecera
\endfirsthead
\cabecera
\endhead
\texttt{codigo} & text & \textbf{PP} \textemdash{} código de la consulta \\
\texttt{producto\_id} & uuid & \textbf{CC} \textemdash{} producto usado \\
cantidad\_usada & int & Cantidad descontada del stock \\
\end{longtable}

% ----- consulta_codigos_by_animal -----
\subsection{\texttt{consulta\_codigos\_by\_animal}}
Índice puntero: consultas de un animal (solo guarda el código, inmutable).

\begin{longtable}{C T R}
\cabecera
\endfirsthead
\cabecera
\endhead
\texttt{animal\_id} & uuid & \textbf{PP} \textemdash{} agrupa por animal \\
\texttt{codigo} & text & \textbf{CC} \textemdash{} código de la consulta \\
\end{longtable}

% ----- consulta_codigos_by_veterinario -----
\subsection{\texttt{consulta\_codigos\_by\_veterinario}}
Índice puntero: consultas de un veterinario.

\begin{longtable}{C T R}
\cabecera
\endfirsthead
\cabecera
\endhead
\texttt{veterinario\_id} & uuid & \textbf{PP} \textemdash{} agrupa por veterinario \\
\texttt{codigo} & text & \textbf{CC} \textemdash{} código de la consulta \\
\end{longtable}

% ============================================================
\section{Notas de modelado}
% ============================================================
\begin{itemize}
  \item \textbf{Desnormalización:} \texttt{animales\_by\_id} y
        \texttt{animales\_by\_rescatista} mantienen los mismos datos del animal
        sincronizados para soportar dos accesos distintos.
  \item \textbf{Tablas puntero:} \texttt{consulta\_codigos\_by\_animal} y
        \texttt{consulta\_codigos\_by\_veterinario} guardan solo el código de la
        consulta (dato inmutable), de modo que no requieren sincronización al
        cambiar el estado o la fecha de una consulta.
  \item \textbf{Borrado lógico:} rescatistas, veterinarios, servicios y productos
        usan la columna \texttt{activo} para darse de baja sin romper referencias
        históricas.
  \item \textbf{Comunicación entre keyspaces:} las relaciones entre servicios
        (p.\,ej. \texttt{animal\_id} o \texttt{veterinario\_id} en
        \texttt{consultas\_by\_codigo}) son referencias por \texttt{uuid}; no hay
        \emph{joins}, cada microservicio es dueño de su propio keyspace.
\end{itemize}

\end{document}
